% 阴影面积的意义——物理量的累积效应
% keys 速度、电荷量、平均物理量
% license Usr
% type Tutor

\subsection{物理量的累积效应}
\subsubsection{速度}
在初中里，我们已经学过“速度”，也就是单位时间内的路程。然而这个概念有很多局限，一方面，路程和时间都是数字，“速度”很自然的，也没有给出其他更多的信息。另一方面，当我们用初中定义去描述任何对象的快慢时，都是默认其快慢不变的。既然小明从西到东的速度是$6 \mathrm{m/s}$，那么在任意的$0.02\mathrm s$内，移动的路程便总是
\subsubsection{冲量}
\subsubsection{电荷量}
\subsection{“平均物理量”}